\begin{proof}[Proof of \cref{obj:prop:symmetry-reduction}]
Write \(n=2m\) and, for a matrix \(X\in\mathbb R^{n\times n}\), introduce the two off-diagonal block sums and the two pair counts
\[
S_{\mathrm{w}}(X)=\sum_{\substack{i\ne j\\ i,j\text{ same community}}}X_{ij},
\qquad
S_{\mathrm{a}}(X)=\sum_{\substack{i,j\\ \text{opposite communities}}}X_{ij},
\]
\[
N_{\mathrm{w}}=2m(m-1),
\qquad
N_{\mathrm{a}}=2m^2 ,
\]
the counts being the numbers of ordered index pairs of each kind.  Call \(G(d,w,c)\) the \emph{block-constant} matrix with \(d\) on the diagonal, \(w\) on within-community off-diagonal entries, and \(c\) on across-community entries.  \cref{obj:ass:two-block-homophily} gives \(m\ge2\), so \(N_{\mathrm{w}},N_{\mathrm{a}}>0\).

\emph{Step 1 (a master trace identity).}  If \(X\) is symmetric, then grouping the \(n^2\) ordered pairs into diagonal, within-community, and across-community pairs gives
\[
\operatorname{tr}\bigl(G(d,w,c)\,X\bigr)
=d\sum_{i}X_{ii}+w\,S_{\mathrm{w}}(X)+c\,S_{\mathrm{a}}(X).
\]
% lean: trace_blockConstMat_mul
The three matrices \(L_m\), \(L_m^\dagger\) and \(J_n\) are block-constant: \(L_m\) has the constant weighted degree \((m-1)a/m+b\) on its diagonal and entries \(-a/m\), \(-b/m\) off it; \(J_n=G(1,1,1)\); and \(L_m^\dagger\) is a linear combination of the identity and the two rank-one projections onto the all-ones and community-sign directions, all three of which are block-constant.  Consequently, if \(X\) and \(Y\) are symmetric with unit diagonal and
\[
S_{\mathrm{w}}(Y)=S_{\mathrm{w}}(X),
\qquad
S_{\mathrm{a}}(Y)=S_{\mathrm{a}}(X),
\]
then \(\operatorname{tr}(L_mY)=\operatorname{tr}(L_mX)\), \(\operatorname{tr}(L_m^\dagger Y)=\operatorname{tr}(L_m^\dagger X)\) and \(\operatorname{tr}(J_nY)=\operatorname{tr}(J_nX)\); if in addition \(\|Y\|_{S_2}\le\|X\|_{S_2}\), then, since \(\kappa\ge0\),
\[
F_{r,\kappa}(Y)\le F_{r,\kappa}(X).
\]
% lean: trace_symmetrize % lean: designObjective_le_of_blockSums
This comparison is the common engine of the two symmetrization claims.

\emph{Step 2 (relaxed matrices: the block average is feasible).}  Let \(X\) be positive semidefinite with \(X_{ii}=1\) for all \(i\), and set
\[
u=\frac{S_{\mathrm{w}}(X)}{N_{\mathrm{w}}},
\qquad
v=\frac{S_{\mathrm{a}}(X)}{N_{\mathrm{a}}},
\]
the averages of the within- and across-community off-diagonal entries.  By construction \(X(u,v)\) is symmetric with unit diagonal and
\[
S_{\mathrm{w}}\bigl(X(u,v)\bigr)=N_{\mathrm{w}}u=S_{\mathrm{w}}(X),
\qquad
S_{\mathrm{a}}\bigl(X(u,v)\bigr)=N_{\mathrm{a}}v=S_{\mathrm{a}}(X).
\]

Membership in \(\mathcal E_m^{\mathrm{blk}}\) requires the three inequalities of \cref{obj:def:block-elliptope}.  First, testing positive semidefiniteness against \(e_i-e_j\) and \(e_i+e_j\) and using the unit diagonal gives
\[
0\le(e_i-e_j)^\top X(e_i-e_j)=2-2X_{ij},
\qquad
0\le(e_i+e_j)^\top X(e_i+e_j)=2+2X_{ij},
\]
so \(|X_{ij}|\le1\) for all \(i,j\). % lean: abs_entry_le_one
Averaging the bound \(X_{ij}\le1\) over the \(N_{\mathrm{w}}\) within-community off-diagonal pairs gives \(u\le1\), i.e.\ \(1-u\ge0\).  Second, with \(s\) the community sign vector (\(+1\) on \(A_m\), \(-1\) on \(B_m\)) the sign product \(s_is_j\) is \(+1\) on diagonal and within-community pairs and \(-1\) on across-community pairs, so
\[
0\le s^\top Xs=\sum_iX_{ii}+S_{\mathrm{w}}(X)-S_{\mathrm{a}}(X)
=2m+2m(m-1)u-2m^2v
=2m\bigl(1+(m-1)u-mv\bigr),
\]
and dividing by \(2m>0\) gives \(1+(m-1)u-mv\ge0\).  Third, testing against the all-ones vector, whose quadratic form is the total entry sum,
\[
0\le\sum_{i,j}X_{ij}=\sum_iX_{ii}+S_{\mathrm{w}}(X)+S_{\mathrm{a}}(X)
=2m\bigl(1+(m-1)u+mv\bigr),
\]
giving \(1+(m-1)u+mv\ge0\).  With \cref{obj:ass:two-block-homophily} this is exactly \(X(u,v)\in\mathcal E_m^{\mathrm{blk}}\). % lean: blockElliptopeMem_symmetrize

\emph{Step 3 (relaxed matrices: the objective does not increase).}  The block sums already match by Step 2, so by Step 1 only the Frobenius term has to be compared.  Squaring entrywise and grouping pairs as before,
\[
\|X\|_{S_2}^2=2m+\sum_{\text{within}}X_{ij}^2+\sum_{\text{across}}X_{ij}^2,
\qquad
\|X(u,v)\|_{S_2}^2=2m+N_{\mathrm{w}}u^2+N_{\mathrm{a}}v^2 .
\]
Cauchy--Schwarz over each of the two pair families gives \(\bigl(\sum_{\text{within}}X_{ij}\bigr)^2\le N_{\mathrm{w}}\sum_{\text{within}}X_{ij}^2\), i.e.
\[
N_{\mathrm{w}}u^2\le\sum_{\text{within}}X_{ij}^2,
\qquad\text{and likewise}\qquad
N_{\mathrm{a}}v^2\le\sum_{\text{across}}X_{ij}^2 .
\]
Hence \(\|X(u,v)\|_{S_2}\le\|X\|_{S_2}\), % lean: frobeniusNorm_symmetrize_le
and Step 1 yields
\[
F_{r,\kappa}\bigl(X(u,v)\bigr)\le F_{r,\kappa}(X).
\]
% lean: symmetrize_objective_le

\emph{Step 4 (assignment laws: the orbit average).}  Let \(H\) be the finite, nonempty group of two-block automorphisms: permutations \(\sigma\) of the units that either preserve both communities setwise or swap them.  For \(P\in\mathcal P_m^{\mathrm{bal}}\) define the orbit average
\[
P^{\mathrm{orb}}(Z=z)=\frac{1}{|H|}\sum_{\sigma\in H}P\bigl(Z=z\circ\sigma\bigr).
\]
It is a probability law, it is sign-symmetric because each summand is, and it is \(H\)-invariant because precomposing \(z\) with \(\tau\in H\) permutes the summation over \(H\); hence \(P^{\mathrm{orb}}\in\mathcal P_m^{\mathrm{sym}}\). % lean: orbitAvgDesign_mem_class
Sign symmetry gives zero one-point margins: replacing \(z\) by \(-z\) leaves the law unchanged and flips the sign of \(Z_i\), so \(\mathbb E[Z_i]=-\mathbb E[Z_i]\) and therefore
\[
\mathbb E[Z_i]=0
\qquad\text{for every }i .
\]
% lean: balanced_zero_margin

\emph{Step 5 (assignment laws: the objective does not increase).}  The same change of variables in the expectation gives the entrywise orbit-averaging identity
\[
X\bigl(P^{\mathrm{orb}}\bigr)_{ij}
=\frac{1}{|H|}\sum_{\sigma\in H}X(P)_{\sigma(i)\sigma(j)} .
\]
% lean: orbitAvg_secondMoment_entry
Every \(\sigma\in H\) maps within-community ordered pairs to within-community ordered pairs and across-community pairs to across-community pairs, and does so bijectively; hence summing the display over each family gives
\[
S_{\mathrm{w}}\bigl(X(P^{\mathrm{orb}})\bigr)=S_{\mathrm{w}}\bigl(X(P)\bigr),
\qquad
S_{\mathrm{a}}\bigl(X(P^{\mathrm{orb}})\bigr)=S_{\mathrm{a}}\bigl(X(P)\bigr).
\]
% lean: Ssame_orbitAvg % lean: Scross_orbitAvg
For the Frobenius term, Cauchy--Schwarz applied to the average over \(H\) gives entrywise
\[
X\bigl(P^{\mathrm{orb}}\bigr)_{ij}^2
\le\frac{1}{|H|}\sum_{\sigma\in H}X(P)_{\sigma(i)\sigma(j)}^2 ,
\]
and summing over all \((i,j)\), each inner sum being a permutation of the full square sum, yields
\[
\bigl\|X(P^{\mathrm{orb}})\bigr\|_{S_2}\le\bigl\|X(P)\bigr\|_{S_2}.
\]
% lean: frobeniusNorm_orbitAvg_le
Both second moments are symmetric with unit diagonal (\(X(P)_{ii}=\mathbb E[Z_i^2]=1\)), so Step 1 applies and gives
\[
F_{r,\kappa}\bigl(X(P^{\mathrm{orb}})\bigr)\le F_{r,\kappa}\bigl(X(P)\bigr).
\]
% lean: design_symmetrize

\emph{Step 6 (an infimum-reduction principle).}  Let \(f\) be a real function on a set \(\mathcal S\), let \(\mathcal T\subseteq\mathcal S\) be nonempty, suppose \(f\) is bounded below on \(\mathcal S\), and suppose every \(x\in\mathcal S\) admits \(y\in\mathcal T\) with \(f(y)\le f(x)\).  Then
\[
\inf_{\mathcal S}f=\inf_{\mathcal T}f .
\]
Indeed \(\inf_{\mathcal S}f\le\inf_{\mathcal T}f\) because \(\mathcal T\subseteq\mathcal S\), and conversely each value \(f(x)\), \(x\in\mathcal S\), dominates some \(f(y)\ge\inf_{\mathcal T}f\). % lean: sInf_image_reduce

\emph{Step 7 (the two infimum identities).}  On the relaxed side take
\[
\mathcal S_{\mathrm{rel}}=\{X:X\succeq0,\ X_{ii}=1\ \forall i\},
\qquad
\mathcal T_{\mathrm{blk}}=\mathcal E_m^{\mathrm{blk}} .
\]
The inclusion \(\mathcal T_{\mathrm{blk}}\subseteq\mathcal S_{\mathrm{rel}}\) follows from the projection decomposition used for block matrices.  Let \(P_{\mathbf 1}=J_n/(2m)\), let \(P_s=ss^\top/(2m)\), and let \(P_\perp=I_n-P_{\mathbf 1}-P_s\).  For every \(u,v\),
\[
X(u,v)=(1-u)P_\perp
+\bigl(1+(m-1)u-mv\bigr)P_s
+\bigl(1+(m-1)u+mv\bigr)P_{\mathbf 1}.
\]
The matrices \(P_{\mathbf 1}\) and \(P_s\) are positive semidefinite rank-one projections.  Also, for any vector \(\xi\), if
\[
A=\sum_{i\in A_m}\xi_i,
\qquad
B=\sum_{i\in B_m}\xi_i,
\qquad
Q_A=\sum_{i\in A_m}\xi_i^2,
\qquad
Q_B=\sum_{i\in B_m}\xi_i^2,
\]
then
\[
\xi^\top P_\perp \xi
=Q_A+Q_B-\frac{A^2+B^2}{m}\ge0
\]
by Cauchy--Schwarz on the two communities, so \(P_\perp\) is positive semidefinite.  Thus the three nonnegative coordinates in \cref{obj:def:block-elliptope} make \(X(u,v)\) positive semidefinite, and the definition of \(X(u,v)\) gives unit diagonal. % lean: blockElliptope_subset_elliptope
The set \(\mathcal T_{\mathrm{blk}}\) is nonempty because \(I_n=X(0,0)\in\mathcal E_m^{\mathrm{blk}}\).

It remains to record the lower bound on the relaxed objective.  Choose block-constant representations
\[
L_m=G(d_L,w_L,c_L),
\qquad
L_m^\dagger=G(d_\dagger,w_\dagger,c_\dagger),
\qquad
J_n=G(1,1,1),
\]
and set
\[
B_L=|d_L|\,2m+|w_L|N_{\mathrm{w}}+|c_L|N_{\mathrm{a}},
\qquad
B_\dagger=|d_\dagger|\,2m+|w_\dagger|N_{\mathrm{w}}+|c_\dagger|N_{\mathrm{a}},
\]
\[
B_J=2m+N_{\mathrm{w}}+N_{\mathrm{a}} .
\]
For \(X\in\mathcal S_{\mathrm{rel}}\), Step 2 gives \(|X_{ij}|\le1\); hence \(|S_{\mathrm{w}}(X)|\le N_{\mathrm{w}}\) and \(|S_{\mathrm{a}}(X)|\le N_{\mathrm{a}}\).  The trace identity in Step 1 then gives, for every block-constant \(G(d,w,c)\),
\[
\bigl|\operatorname{tr}(G(d,w,c)X)\bigr|
\le |d|\,2m+|w|N_{\mathrm{w}}+|c|N_{\mathrm{a}} .
\]
Therefore
\[
F_{r,\kappa}(X)
\ge -B_L-|r|B_\dagger+\kappa\|X\|_{S_2}-B_J
\ge -\bigl(B_L+|r|B_\dagger+B_J\bigr),
\]
because \(\kappa\ge0\) and \(\|X\|_{S_2}\ge0\). % lean: designObjective_bddBelow
Steps 2--3 supply, for every \(X\in\mathcal S_{\mathrm{rel}}\), the point \(X(u,v)\in\mathcal T_{\mathrm{blk}}\) with no larger objective.  Step 6 gives
\[
\inf_{\substack{X\succeq0\\ X_{ii}=1\ \forall i}}F_{r,\kappa}(X)
=\inf_{X\in\mathcal E_m^{\mathrm{blk}}}F_{r,\kappa}(X).
\]

On the implementable side take
\[
\mathcal S_{\mathrm{bal}}=\mathcal P_m^{\mathrm{bal}},
\qquad
\mathcal T_{\mathrm{sym}}=\mathcal P_m^{\mathrm{sym}},
\qquad
f(P)=F_{r,\kappa}\bigl(X(P)\bigr).
\]
Then \(\mathcal T_{\mathrm{sym}}\subseteq\mathcal S_{\mathrm{bal}}\) by \cref{obj:def:block-exchangeable-design-class}; \(\mathcal T_{\mathrm{sym}}\ne\emptyset\) because the iid fair-sign law is block-exchangeable; % lean: iidDesign_mem_blockExchangeable
and \(f\) is bounded below because every second moment \(X(P)=\mathbb E_P[ZZ^\top]\) is positive semidefinite with unit diagonal, hence belongs to \(\mathcal S_{\mathrm{rel}}\) and satisfies the lower bound just displayed. % lean: assignmentSecondMoment_posSemidef
Steps 4--5 supply, for every \(P\in\mathcal S_{\mathrm{bal}}\), the law \(P^{\mathrm{orb}}\in\mathcal T_{\mathrm{sym}}\) with no larger objective.  Step 6 gives
\[
\inf_{P\in\mathcal P_m^{\mathrm{bal}}}F_{r,\kappa}\bigl(X(P)\bigr)
=\inf_{P\in\mathcal P_m^{\mathrm{sym}}}F_{r,\kappa}\bigl(X(P)\bigr).
\]
% lean: symmetry_reduction
\end{proof}
